\documentclass[11pt]{article}
\usepackage[margin=1in]{geometry}
\usepackage{amsmath,amssymb,amsthm}
\usepackage{graphicx}
\usepackage{hyperref}
\usepackage{microtype}

\newtheorem{theorem}{Theorem}
\newtheorem{lemma}{Lemma}
\newcommand{\Aut}{\operatorname{Aut}}
\newcommand{\Homeo}{\operatorname{Homeo}}

\title{Finite-tree automorphism actions are Tame$_1$ but not Tame$_2$\\
\large A candidate partial answer to Question 5.12 of arXiv:2601.01681}
\author{Automated mathematics run \texttt{fa\_banach\_001}, lane 3}
\date{9 August 2026}

\begin{document}
\maketitle

\noindent\textbf{Status.} Candidate partial result; likely valid; awaiting human review.

\section{Source question and scope}

Megrelishvili asks which natural finite-rank compact metrizable median
$G$-algebras are Tame$_1$ or Tame$_2$ \cite[Question 5.12]{Megrelishvili}.
The source gives the interval and finite-dimensional cubes as Tame$_1$
examples, and contrasts them with the Wa\.{z}ewski dendrite, whose full
automorphism action is not Tame$_1$.

\begin{center}
\includegraphics[width=\textwidth]{figures/open_problem_crop.png}
\end{center}

\noindent\textit{Transcription.} ``Which natural finite rank compact metrizable
median $G$-algebras $K$ are Tame$_1$? Tame$_2$?''

This packet gives a complete classification for the natural rank-one class of
geometric realizations of finite simplicial trees. It does not answer the
source question for arbitrary finite-rank median algebras or arbitrary group
actions.

\section{Definitions}

Let $K$ be the geometric realization of a finite simplicial tree. Its median
$m(x,y,z)$ is the unique point in
\[
[x,y]\cap[y,z]\cap[z,x],
\]
where $[x,y]$ denotes the unique arc from $x$ to $y$. We give
$G=\Aut(K)$ the compact-open topology and let it act naturally on $K$.

For a compact $G$-space $K$, its enveloping (Ellis) semigroup is
\[
E(G,K)=\overline{\{g:K\to K:g\in G\}}^{\,K^K},
\]
where the closure is pointwise. In the terminology used by the source,
$(G,K)$ is Tame$_1$ when $E(G,K)$ is first countable, and Tame$_2$ when
$E(G,K)$ is hereditarily separable (within the tame class).

The source recalls that the enveloping semigroup of the order-preserving
homeomorphism action on $[0,1]$ is a compact subspace of the first-countable
Helly space of monotone maps, but is not hereditarily separable
\cite[Example 5.13]{Megrelishvili}; see also \cite{GlasnerMegrelishvili}.

\section{Proof intuition}

A finite tree has a finite topological skeleton: its endpoints and branch
points. Removing this skeleton leaves finitely many open edges. Every
homeomorphism permutes the skeleton and the edges. After fixing this finite
permutation, an automorphism is described on each edge by a monotone interval
homeomorphism. Pointwise limits are therefore controlled by finitely many
Helly compacta, which gives first countability.

The negative Tame$_2$ conclusion comes from looking in the opposite direction:
one edge already carries all order-preserving interval homeomorphisms, extended
by the identity on the rest of the tree. Their non-hereditarily-separable Ellis
semigroup embeds in the Ellis semigroup of the full tree action.

\section{The partial result}

\begin{theorem}\label{thm:main}
Let $K$ be the geometric realization of a nondegenerate finite simplicial tree,
equipped with its canonical median, and let
$G=\Aut(K)$ be the compact-open group of all homeomorphic median
automorphisms. Then the natural compact median $G$-algebra $(G,K)$ is
Tame$_1$ but not Tame$_2$.
\end{theorem}

\begin{proof}
Every homeomorphism of a tree preserves unique arcs and therefore preserves
their triple-intersection median. Thus $G=\Homeo(K)$ as sets.

Let $V$ be the set of endpoints and branch points of $K$. It is finite and is
preserved by every homeomorphism. The connected components of $K\setminus V$
are finitely many open arcs; their closures will be called the essential
edges. Every $g\in G$ permutes $V$ and the essential edges.

Put $E=E(G,K)$. If $p\in E$, choose a net $(g_i)$ in $G$ converging pointwise
to $p$. For each $v\in V$, the net $(g_i v)$ takes values in the finite set
$V$. Its convergence forces it to be eventually constant. Since $V$ is
finite, on one common tail all $g_i|_V$ equal a single permutation $\sigma$.
Consequently every $p\in E$ belongs to
\[
E_\sigma:=\overline{\{g\in G:g|_V=\sigma\}}^{\,K^K}
\]
for one of finitely many realizable permutations $\sigma$.

These pieces are pairwise disjoint and clopen in $E$: evaluation on the finite
invariant set $V$ separates the possible restrictions. Fix $\sigma$. It also
fixes the target essential edge $\sigma(e)$ of each essential edge $e$. Choose
coordinates $\theta_e:[0,1]\to e$. After reversing one coordinate when the
endpoint order demands it, every normalized restriction
\[
\theta_{\sigma(e)}^{-1}\circ g|_e\circ\theta_e
\]
is an increasing homeomorphism of $[0,1]$. Hence restriction to the finitely
many edges defines a continuous injection of $E_\sigma$ into a finite product
of Helly compacta of nondecreasing interval maps. Since $E_\sigma$ is compact
and the product is Hausdorff, this injection is a topological embedding.
The Helly compactum is first countable \cite[Example 5.13]{Megrelishvili}; so
finite products and their subspaces are first countable. Thus every
$E_\sigma$ is first countable. Because the finitely many $E_\sigma$ are clopen,
$E$ is first countable. Therefore $(G,K)$ is Tame$_1$.

It remains to rule out Tame$_2$. Fix one essential edge $e$ and identify it
with $[0,1]$. Every increasing homeomorphism of $e$ fixing its endpoints
extends by the identity on $K\setminus e$ to an element of $G$. Let $H$ be this
subgroup. Extension by the identity induces a topological embedding
\[
E(H_+([0,1]),[0,1])\hookrightarrow E(H,K)\subseteq E(G,K).
\]
The interval enveloping semigroup on the left is not hereditarily separable
\cite[Example 5.13]{Megrelishvili}. Hence $E(G,K)$ is not hereditarily
separable. The system is not Tame$_2$.
\end{proof}

\section{Verification notes}

The proof deliberately uses the fact that the skeleton-permutation pieces are
\emph{clopen}. An arbitrary finite union of closed first-countable subspaces
need not be first countable, so merely writing the Ellis semigroup as such a
union would not suffice. Evaluation on the finite invariant skeleton supplies
the needed clopen decomposition.

The interval case is recovered when $V$ consists of two endpoints and there is
one essential edge. Adding or deleting valence-two subdivision vertices does
not change $V$, the essential edges, or the argument. A triod is the first
branched example: the center and three endpoints give a four-point skeleton
and three Helly coordinates in each permutation component.

No computational step is used. The proof depends externally only on the
first-countability and non-hereditary-separability facts for the classical
interval Ellis semigroup, both explicitly recorded in the source's
Example 5.13.

\section{Limitations and novelty check}

This is a partial answer: it treats the full automorphism action on finite
topological trees, not arbitrary actions on dendrites and not higher-rank
median spaces. It also does not classify which infinite dendrites have
Tame$_1$ full automorphism actions.

On 9 August 2026, the run's lightweight registry, solution, attempt, and
proof-gap indexes were searched for arXiv:2601.01681 and finite-rank
median/tameness terms, with no hit. A bounded web/arXiv search used combinations
of ``finite tree'', ``automorphism'', ``Tame$_1$'', ``dendrite'', and
``enveloping semigroup''. It found the source paper and Codenotti's dendrite
counterexamples \cite{Codenotti}, but no prior finite-tree classification.
Novelty confidence is therefore moderate, not definitive.

\section{Human-review recommendation}

Send to a specialist in tame dynamics or continuum theory. The highest-value
checks are: (i) that the essential-skeleton restriction indeed yields clopen
Ellis components, and (ii) that extension by the identity embeds the complete
interval Ellis semigroup, not merely the acting group.

\begin{thebibliography}{9}

\bibitem{Megrelishvili}
M. Megrelishvili,
\emph{Tameness of actions on finite rank median algebras},
arXiv:2601.01681, April 2026.

\bibitem{GlasnerMegrelishvili}
E. Glasner and M. Megrelishvili,
\emph{Todor\v{c}evi\'c's Trichotomy and a hierarchy in the class of tame
dynamical systems},
Trans. Amer. Math. Soc. \textbf{375} (2022), 4513--4548.

\bibitem{Codenotti}
A. Codenotti,
\emph{Some examples of tame dynamical systems answering questions of Glasner
and Megrelishvili},
Proc. Amer. Math. Soc. \textbf{153} (2025), 2433--2449;
arXiv:2305.00726.

\end{thebibliography}

\end{document}
